\section{Nilpotent quotients and bounded projections}
\label{sec:structure}

The remaining growth case has a polynomial bound along a sequence of
scales. Its useful consequence is an action on the original graph with
a nilpotent quotient. The quotient will supply coordinates; all
percolation paths and edge exposures will remain in the original graph.
We prove the following action theorem and then apply it to the action
constructed from the polynomial scales.

\begin{theorem}[Criticality with a nilpotent quotient action]
\label{thm:structured}
Let $G$ be an infinite connected locally finite graph of subexponential
growth. Suppose a group $\Lambda$ acts by automorphisms on $G$ with
finitely many vertex orbits and admits an epimorphism $\varrho:\Lambda\to N$ onto a
finitely generated nilpotent group, such that $\ker\varrho$ has finite
vertex orbits and $\varrho(\Lambda_v)$ is finite for every vertex $v$.
Here $\Lambda_v=\{g\in\Lambda:gv=v\}$ is the vertex stabilizer.
If $\pc(G)<1$, then
\[
 \PP_{\pc(G)}(|C_v|=\infty)=0\qquad(v\in G).
\]
Consequently, critical percolation has almost surely no infinite cluster.
\end{theorem}

The action need not be faithful, and \(\Lambda\) need not itself be
nilpotent. We first show that a polynomial bound along an unbounded
sequence supplies exactly these action data. We then prove the theorem
by induction on a cyclic series of \(N\). In the induction step, a
component over a bounded coordinate box has a quotient with fewer
infinite cyclic factors. This excludes rays with bounded projection;
Sections~\ref{sec:moves} and~\ref{sec:corridor-exploration} use the
remaining rays to construct percolation below \(\pc\).

\subsection{From polynomial scales to the quotient action}

The structural ancestry is Gromov's group theorem \cite{Gromov1981},
Trofimov's graph and near-polynomial growth theorems
\cite{Trofimov1985,Trofimov2003}, and Tessera and Tointon's one-scale
form \cite{TesseraTointon2021}. The proof below uses the latter with
its action and stabilizer conclusions.

\begin{proposition}[The action supplied by a polynomial scale]
\label{prop:nilpotent-model}
Let $G$ be an infinite connected locally finite quasi-transitive graph.
Suppose that for some $d<\infty$ there are radii $R_j\to\infty$ such that
$V_-(R_j)\le R_j^d$. Then there are a finite-index subgroup
$\Lambda\le\operatorname{Aut}(G)$, a finitely generated nilpotent group
$N$, and an epimorphism
\[
 \varrho:\Lambda\longrightarrow N
\]
with the following properties:
\begin{enumerate}
 \item $\Lambda$ has finitely many vertex orbits;
 \item $\ker\varrho$ has finite vertex orbits;
 \item $\varrho(\Lambda_v)$ is finite for every vertex $v$, where
 $\Lambda_v=\{g\in\Lambda:gv=v\}$.
\end{enumerate}
\end{proposition}

\begin{proof}
Increase $d$ if necessary so that $d\ge1$. Let $C$ be a common density
radius for the finitely many vertex orbits of $\operatorname{Aut}(G)$. Fix one orbit $O$, set $L=2C+1$, and form a
graph $\Gamma$ on $O$ by joining distinct vertices whose distance in $G$
is at most $L$. To connect two points of $O$, take a path between them
in $G$ and replace its intermediate vertices by points of $O$ within
distance $C$. Consecutive replacements are equal or adjacent in
$\Gamma$. Thus $\Gamma$ is connected. It is locally finite, and the
image $A$ of $\operatorname{Aut}(G)$ on $O$ is a transitive subgroup of
$\operatorname{Aut}(\Gamma)$.

At a radius $R=R_j$, choose $x$ with $|B(x,R)|=V_-(R)$ and $u\in O$
with $d(x,u)\le C$. For $n=\lfloor(R-C)/L\rfloor$,
\[
 B_\Gamma(u,n)\subseteq B_G(u,Ln)\subseteq B_G(x,R).
\]
Writing $\beta_\Gamma(n)$ for the cardinality of a radius-$n$ ball in
this transitive graph, we obtain $\beta_\Gamma(n)\le R^d$.
For all sufficiently large such scales, $R\le c n$ with a fixed
constant $c$. Absorbing $c^d$ into an extra power of $n$ gives
\[
 \beta_\Gamma(n)\le n^{d+1}\beta_\Gamma(1)
\]
for arbitrarily large integers $n$.

Apply Tessera--Tointon's Corollary~2.4 with growth exponent $d+1$,
$\lambda=1/2$, and the transitive group $A$
\cite[Corollary~2.4]{TesseraTointon2021}.
Its hypotheses require only one sufficiently large scale and do not
require $A$ to be closed in $\operatorname{Aut}(\Gamma)$.
The conclusion supplies a normal subgroup $H\lhd A$ whose orbits
have finite diameter, such that $Q=A/H$ acts faithfully on
$\Gamma/H$, is finitely generated and virtually nilpotent, and has
finite vertex stabilizers. The $H$-orbits are finite because $\Gamma$
is locally finite.

Let $\vartheta:\operatorname{Aut}(G)\to Q$ be the composite
epimorphism, and put $K=\ker\vartheta$. Its orbits on $O$ are the
$H$-orbits. For an arbitrary $v\in G$, choose $u\in O$ with
$d(u,v)\le C$. Then
\[
 Kv\subseteq\bigcup_{w\in Ku}B_G(w,C),
\]
so $Kv$ is finite. Also, $\operatorname{Aut}(G)_v$ moves $u$ within
the finite ball $B_G(v,C)$. Hence its image under $\vartheta$ has a
finite orbit on the block $H(u)\in\Gamma/H$. Its stabilizer at that
block is a subgroup of the finite $Q$-stabilizer. The orbit--stabilizer
formula therefore makes $\vartheta(\operatorname{Aut}(G)_v)$ finite.

Choose a finite-index nilpotent subgroup $N\le Q$ and let
$\Lambda=\vartheta^{-1}(N)$. A finite-index subgroup of a finitely
generated group is finitely generated: rewriting a word through a
finite set of coset representatives gives a finite generating set.
Thus $N$ is finitely generated. Each $\operatorname{Aut}(G)$-orbit
splits into at most $[\operatorname{Aut}(G):\Lambda]$ many
$\Lambda$-orbits. Restricting $\vartheta$ to $\Lambda$ gives
$\varrho$, with the other two properties already proved.
\end{proof}

\subsection{Cyclic series and the decreasing parameter}

A \emph{cyclic series} is a finite subnormal series
\[
 N=N_0\triangleright N_1\triangleright\cdots\triangleright N_k=\{1\}
\]
whose factors $N_i/N_{i+1}$ are cyclic, possibly finite or trivial.
We use the number of its infinite factors. The induction applies to
all data equipped with a chosen series of at most a given count;
independence of that count from the choice of series is not needed.

The group facts below are classical; see
\cite[Theorems~5.2.5 and~5.2.18]{Robinson1996}.
We give the arguments to make the decrease of the induction parameter
explicit.

\begin{lemma}[Cyclic series and their subgroups]
\label{lem:rank}
Every finitely generated nilpotent group $N$ has a cyclic series.
If such a series has $h$ infinite factors, then every subgroup $J\le N$
is finitely generated and inherits a cyclic series with at most $h$
infinite factors. If $[N:J]=\infty$, the inherited count is strictly
less than $h$.
Moreover, if the finitely generated abelianization $N/[N,N]$ has rank
at most one, then $[N,N]$ is finite. In this case $N$ is either finite
or admits an epimorphism onto $\mathbb Z$ with finite kernel.
\end{lemma}

\begin{proof}
Write $\gamma_1=N$ and $\gamma_{i+1}=[N,\gamma_i]$ for the lower
central series. The first factor $N/\gamma_2$ is finitely generated
abelian. Modulo $\gamma_{i+2}$, the commutator product identities give
a bilinear map
\[
 (N/\gamma_2)\times(\gamma_i/\gamma_{i+1})
 \longrightarrow\gamma_{i+1}/\gamma_{i+2},
 \qquad (x\gamma_2,y\gamma_{i+1})\longmapsto[x,y]\gamma_{i+2},
\]
whose values generate the target. Here bilinearity follows because
commutators of one higher weight vanish in the target quotient.
Inductively, images of pairs of finite generating sets generate each
successive factor. These factors are abelian, and nilpotence makes
only finitely many of them nontrivial. Refining each finitely generated
abelian factor by adjoining its generators one at a time gives a
cyclic series.

For a subgroup $J$, intersect this series with $J$. The natural map
\[
 (J\cap N_i)/(J\cap N_{i+1})\longrightarrow N_i/N_{i+1}
\]
is injective, so each inherited factor is cyclic. Lifting generators
successively up this finite series shows that $J$ is finitely generated.
Suppose its count of infinite factors is still $h$. Then its image in
every infinite cyclic ambient factor is infinite, hence has finite
index; its image in a finite factor also has finite index. Induction
up the series gives $[N:J]<\infty$. Explicitly, if
$J_i=J\cap N_i$, then
\[
 [J_iN_{i+1}:J_i]=[N_{i+1}:J\cap N_{i+1}],
\]
and $[N_i:J_iN_{i+1}]$ is the index of the image in the cyclic factor.
Both indices are finite at each step, starting from the trivial last
group. This proves the strict decrease for infinite-index subgroups.

For the last assertion, the first commutator pairing is alternating
on $N/\gamma_2$. If that abelian group has rank at most one, the
pairing has only torsion values: its exterior square is a finitely
generated torsion group. Its finitely many generating values make
$\gamma_2/\gamma_3$ finite. The displayed bilinear pairing then makes every later factor
finite as well. Thus $[N,N]=\gamma_2$ is finite. Rank zero of the
abelianization makes $N$ finite. In rank one, projection of the
abelianization onto its infinite cyclic factor gives an epimorphism
$N\to\mathbb Z$ with finite kernel.
\end{proof}

\subsection{The induction and equivariant coordinates}

Fix action data \((G,\Lambda,\varrho,N)\) as in
Theorem~\ref{thm:structured}, together with a cyclic series having
\(h\) infinite factors. We prove the theorem simultaneously for all
such data by induction on \(h\). A series with no infinite factor
makes \(N\) finite. Each $\Lambda$-orbit is a finite union of
$\ker\varrho$-orbits and is therefore finite. There are finitely
many $\Lambda$-orbits, contrary to infinitude of $G$.
This is the base case. For the induction step, assume the theorem
for every such action whose quotient has a cyclic series with fewer
than \(h\) infinite factors. Suppose, towards a contradiction, that
\(\pc(G)<1\) and critical percolation has an infinite cluster with
positive probability.

We first construct the coordinates and dispose of the one-dimensional
case. Suppose $\phi:N\to\mathbb Z^j$ is an epimorphism, where $j=1$ or $2$.
Choose representatives $v_1,\ldots,v_k$ for the $\Lambda$-orbits and
put
\begin{equation}\label{eq:equivariant-projection}
 \pi(gv_i)=\phi\varrho(g),\qquad g\in\Lambda.
\end{equation}
This is well defined: if $gv_i=g'v_i$, then $(g')^{-1}g$ stabilizes
$v_i$, its $\varrho$-image belongs to a finite group, and $\phi$
annihilates that image because $\mathbb Z^j$ is torsion free. Moreover,
\begin{equation}\label{eq:projection-equivariance}
 \pi(gx)=\phi\varrho(g)+\pi(x).
\end{equation}
The map is onto, even on each vertex orbit. There are finitely many
oriented edge orbits, so equivariance gives a finite bound
\begin{equation}\label{eq:projection-edge-bound}
 D_\pi:=\sup_{x\sim y}|\pi(x)-\pi(y)|<\infty.
\end{equation}
The norm is Euclidean when $j=2$.

Suppose that $\phi:N\to\mathbb Z$ has finite kernel.
Then $L_0=\ker(\phi\varrho)$ has finite orbits: its image under
$\varrho$ is finite, and its kernel has finite orbits. Thus
\[
 \pi^{-1}(0)=\bigcup_{i=1}^k L_0v_i
\]
is finite. An element with $\phi\varrho$-value $t$ bijects this level
onto $\pi^{-1}(t)$, so all integer levels have the same finite size.

Choose an integer $J\ge\max(1,D_\pi)$. Edges crossing the cut from
$\{\pi\le t\}$ to $\{\pi>t\}$ have endpoints in a fixed number of
levels around $t$. Their number is uniformly bounded, since degree
and level sizes are bounded. The crossing-edge sets for levels
$t\in(J+1)\mathbb Z$ are disjoint. For every $p<1$, their all-closed
events are independent with probabilities bounded below by a positive
constant. Almost surely infinitely many such cuts are closed in both
directions. Every open cluster is then confined between two closed
cuts and lies in finitely many finite levels. This proves
\begin{equation}\label{eq:cyclic-critical-one}
 \pc(G)=1
 \quad\text{whenever $N$ has an epimorphism to $\mathbb Z$
 with finite kernel.}
\end{equation}

Consequently, the hypotheses of Theorem~\ref{thm:structured} with
$\pc(G)<1$ force the abelianization of $N$ to have rank at least two.
It therefore admits an epimorphism onto $\mathbb Z^2$.

\subsection{Excluding rays with bounded projection}

We now choose an epimorphism \(\phi:N\to\mathbb Z^2\), as justified
above, and let \(\pi\) be the associated coordinate map. The smaller-count
induction applies to a component over a bounded box because its
preserving subgroup maps into \(\ker\phi\), an infinite-index subgroup
of \(N\). The next lemma checks all the required action properties.

\begin{lemma}[Exclusion of bounded projections]
\label{lem:bounded-projection}
\label{lem:unbounded-projection}
Let $(G,\Lambda,\varrho,N)$ satisfy the hypotheses of
Theorem~\ref{thm:structured}, with $\pc(G)<1$. Fix a cyclic series for
$N$ with $h$ infinite factors and an epimorphism $\phi:N\to\mathbb Z^2$.
Let $\pi$ be given by \eqref{eq:equivariant-projection}. Assume
Theorem~\ref{thm:structured} for every such action whose nilpotent
quotient has a cyclic series with fewer than $h$ infinite factors.
Then critical percolation on $G$ almost surely has no infinite open ray
whose $\pi$-projection is bounded.
\end{lemma}

\begin{proof}
Suppose otherwise. There are countably many finite boxes in
$\mathbb Z^2$ and countably many connected components of their inverse
images. Hence some fixed finite box $Q$ and a fixed infinite connected
component $G'$ of the induced graph $G[\pi^{-1}(Q)]$ have an infinite
open cluster at $p=\pc(G)$ with positive probability. Subgraph
monotonicity gives $\pc(G')\ge\pc(G)$; this positive probability
gives the reverse inequality. Thus
\[
 \pc(G')=\pc(G)<1,
\]
and $G'$ percolates at its own critical value.

Let $\Lambda'$ consist of the elements of $\ker(\phi\varrho)$ that
preserve $G'$. It has finitely many vertex orbits on $G'$. Indeed,
if $x,y\in G'$ lie in the same $\Lambda$-orbit and satisfy
$\pi(x)=\pi(y)$, choose $g\in\Lambda$ with $gx=y$.
Equation~\eqref{eq:projection-equivariance} gives
$\phi\varrho(g)=0$. This element preserves the induced graph
$G[\pi^{-1}(Q)]$ and maps the component of $x$ to that of $y$;
these are the same component. Thus $g\in\Lambda'$. There are only
finitely many pairs of an original $\Lambda$-orbit and a level in $Q$,
which proves the claim.

Set $N'=\varrho(\Lambda')$. Restriction gives an epimorphism
$\Lambda'\to N'$ whose kernel has finite orbits and whose vertex
stabilizers have finite images. The group $N'$ is nilpotent and is
finitely generated by Lemma~\ref{lem:rank}. It lies in $\ker\phi$,
which has infinite index in $N$, so its inherited cyclic series has
strictly fewer than $h$ infinite factors. The graph $G'$ is connected,
locally finite and infinite. Its intrinsic balls are contained in
ambient balls of the same radius, so it has subexponential growth.
It therefore satisfies every hypothesis of the smaller-count
induction. Theorem~\ref{thm:structured} at that count forbids the
critical infinite cluster just found, a contradiction.
\end{proof}

Thus the induction step has onto equivariant coordinates in
\(\mathbb Z^2\), bounded coordinate jumps, and no infinite open ray
trapped over a bounded set. The following two sections construct an
infinite cluster at a parameter below \(\pc(G)\) from the assumed
critical jump. That contradiction completes the induction and proves
Theorem~\ref{thm:structured}. Proposition~\ref{prop:nilpotent-model}
then applies it to the remaining growth case.
